\chapter{Recomendaciones y Conclusiones}

\section{Conclusiones Principales}

\subsection{Análisis de Flujo de Carga}

Ante la entrada en servicio de la generación en el nodo 3272966 del circuito 10-502:

\begin{itemize}
    \item La cargabilidad de líneas y del transformador de conexión se mantiene dentro de los límites de operación.
    \item La variación máxima de tensión es de 0,00062~p.u.
    \item No aumenta la cargabilidad de las líneas evaluadas: la energía se consume localmente y se alivia el alimentador.
\end{itemize}

\subsection{Análisis de Cortocircuito}

\begin{itemize}
    \item El aporte del SSFV en MT es del orden de 0,008~kA frente a $\approx$5,12~kA trifásicos en el POC.
    \item La variación de $I_k$ es inferior al 0,01\% respecto al caso base.
    \item No se requieren cambios en el relé MICOM P142 ni en los ajustes 50/51/50N/51N del circuito.
\end{itemize}

\subsection{Protecciones y Pérdidas}

\begin{itemize}
    \item Se cumple el Acuerdo CNO 1322 (27, 59, 81 U/O, anti-isla).
    \item Las pérdidas técnicas del sistema se reducen hasta 0,34~kW en las horas de análisis.
\end{itemize}

\section{Recomendaciones}

\subsection{Medición de Calidad de Potencia}

Se recomienda realizar mediciones de calidad de potencia una vez entre en servicio el proyecto, para verificar que el factor de potencia se mantenga dentro del rango regulatorio. Si el FP baja por reducción de potencia activa con reactiva constante, se podrá valorar compensación capacitiva.

\subsection{Uso de Resultados de Cortocircuito}

Los niveles de falla calculados deben usarse para:

\begin{enumerate}
    \item Especificar capacidades de interrupción de breakers, fusibles y seccionadores.
    \item Evaluar saturación de transformadores de corriente.
    \item Diseñar/verificar mallas de puesta a tierra.
    \item Especificar nuevos equipos en ampliaciones futuras.
\end{enumerate}

\subsection{Validación en Campo}

Antes y después de la puesta en servicio conviene medir:

\begin{itemize}
    \item Tensiones en el POC y nodos aledaños
    \item Cargabilidad del transformador de 150~kVA
    \item Calidad de la potencia inyectada (armónicos, FP, flicker)
\end{itemize}

\subsection{Coordinación con ESSA}

Mantener coordinación permanente con el OR para reportar anomalías y actualizar ajustes si cambian las condiciones del circuito 10-502.

\section{Hallazgos Positivos}

\begin{enumerate}
    \item Impacto mínimo en tensiones y cortocircuito.
    \item Transformador de conexión con holgura operativa ($<$56\%).
    \item Reducción de pérdidas técnicas del alimentador.
    \item Consumo local predominante de la generación.
    \item Protecciones existentes del circuito adecuadas y no afectadas.
\end{enumerate}

\section{Concepto Técnico}

El presente estudio de conexión simplificado concluye que:

\begin{quote}
\textit{``El proyecto de generación solar fotovoltaica de 141,6~kWp propuesto para el Supermercado MAS X MENOS sede Guarín puede conectarse al Sistema de Distribución Local de ESSA en el nodo 3272966 del circuito 10-502 (SE Conucos, 13,8~kV) sin generar violaciones a los parámetros de operación normal de la red. La tecnología de inversores limita efectivamente el aporte de cortocircuito. Las protecciones existentes cumplen los requisitos regulatorios y no requieren modificaciones. La instalación contribuye a la reducción de pérdidas técnicas en la red local.''}
\end{quote}

Por lo anterior, se emite \textbf{concepto técnico favorable} a la conexión del proyecto, recomendando las medidas de monitoreo y validación de calidad de potencia descritas en este documento.
